
\documentclass{article}
\usepackage{colortbl}
\usepackage{makecell}
\usepackage{multirow}
\usepackage{supertabular}

\begin{document}

\newcounter{utterance}

\centering \large Interaction Transcript for game `bbh', experiment `snarks', episode 26 with reflexion{-}Qwen3.5{-}27B.
\vspace{24pt}

{ \footnotesize  \setcounter{utterance}{1}
\setlength{\tabcolsep}{0pt}
\begin{supertabular}{c@{$\;$}|p{.15\linewidth}@{}p{.15\linewidth}p{.15\linewidth}p{.15\linewidth}p{.15\linewidth}p{.15\linewidth}}
    \# & \multicolumn{2}{c}{Player} && \multicolumn{2}{c}{Game Master} \\
    \hline

    \theutterance \stepcounter{utterance}  
    & & & \multicolumn{4}{p{0.6\linewidth}}{
        \cellcolor[rgb]{0.9,0.9,0.9}{
            \makecell[{{p{\linewidth}}}]{
                \texttt{\tiny{[P1$\langle$GM]}}
                \texttt{You have attempted this same task before and did not succeed. These are the plans you wrote after those attempts.} \\
\texttt{Trial \#1: New plan: My previous attempt failed because I selected option (B) based on the irony of calling trolls "endangered," but the correct answer was likely (A). I misidentified the sarcastic target; in option (A), the phrase "feed the elephants" is a nonsensical or absurd substitution for the common idiom "feed the trolls," creating a humorous, ironic contrast by treating a mundane or impossible action (feeding elephants) as a standard social interaction, or perhaps the sarcasm lies in the absurdity of the statement itself compared to the expected idiom. Alternatively, I may have misjudged the "endangered" claim as the primary sarcastic element when the sarcasm in (A) relies on the deliberate misuse of the idiom to mock the situation. Next time, I will prioritize identifying which option contains a deliberate, absurd deviation from a common idiom or a statement that is clearly false in a way that mocks the subject, rather than assuming the option with the most obvious "ironic" adjective (like "endangered" for trolls) is the answer. I will specifically look for the option where the literal meaning is so absurd or contradictory to reality that it serves as a clear critique, even if the other option also contains irony. In this specific case, I should have recognized that "feed the elephants" is the sarcastic twist on "feed the trolls," making (A) the intended answer, as the absurdity of the substitution is the core of the humor.} \\
\\ 
\\ 
\texttt{Determine which of two sentences is sarcastic.} \\
\\ 
\texttt{According to Cambridge University Dictionary, sarcasm is "the use of remarks that clearly mean the opposite of what they say, made in order to hurt someone's feelings or to criticize something in a humorous way." Sarcastic sentences often contain satirical or ironic utterances, hyperboles, ambivalent or witty remarks.} \\
\\ 
\texttt{Q: Which statement is sarcastic?} \\
\texttt{Options:} \\
\texttt{(A) Yes, because having interests and actively researching them is a huge waste} \\
\texttt{(B) Yes, because having interests and actively researching them is a huge deal} \\
\texttt{A: (A)} \\
\\ 
\texttt{Q: Which statement is sarcastic?} \\
\texttt{Options:} \\
\texttt{(A) No one is going to disagree with you on this. Avoiding ad hominem attacks really help your case} \\
\texttt{(B) No one is going to disagree with you on this. Ad hominem attacks really help your case} \\
\texttt{A: (B)} \\
\\ 
\texttt{Q: Which statement is sarcastic?} \\
\texttt{Options:} \\
\texttt{(A) Consistency in the league's punishments? What do you think this is supposed to be, politics?} \\
\texttt{(B) Consistency in the league's punishments? What do you think this is supposed to be, moral?} \\
\texttt{A: (A)} \\
\\ 
\texttt{Q: Which statement is sarcastic?} \\
\texttt{Options:} \\
\texttt{(A) You're welcome my dear friend. Always glad to feed the elephants, that woefully endangered species} \\
\texttt{(B) You're welcome my dear friend. Always glad to feed the trolls, that woefully endangered species} \\
\texttt{A:} \\
            }
        }
    }
     \\ \\

    \theutterance \stepcounter{utterance}  
    & \multicolumn{4}{p{0.6\linewidth}}{
        \cellcolor[rgb]{0.561,0.737,0.561}{
            \makecell[{{p{\linewidth}}}]{
                \texttt{\tiny{[P1$\rangle$GM]}}
                \texttt{Based on the definition of sarcasm provided and the pattern of the previous examples, here is the analysis:} \\
\\ 
\texttt{1.  **Analyze the Idiom**: The common internet idiom is "feed the trolls," which means to engage with people who post inflammatory comments to provoke a reaction.} \\
\texttt{2.  **Analyze Option (B)**: "Always glad to feed the trolls, that woefully endangered species." This sentence uses the correct idiom but adds the phrase "woefully endangered species." Trolls are generally considered a nuisance that is abundant, not endangered. Calling them "endangered" is ironic and sarcastic, mocking the idea that they are rare or in need of protection.} \\
\texttt{3.  **Analyze Option (A)**: "Always glad to feed the elephants, that woefully endangered species." This sentence replaces "trolls" with "elephants." While elephants can be endangered, the phrase "feed the elephants" is not the standard idiom for online interaction. This makes the sentence sound like a literal (albeit odd) statement about animals or a complete non{-}sequitur, rather than a sharp, ironic twist on a specific social behavior.} \\
\texttt{4.  **Re{-}evaluating based on the "Trial \#1" context**: The user's previous attempt failed because they chose (B) thinking it was the answer, but the system indicated (A) was correct (or the user *thought* (A) was correct in their reflection). However, looking strictly at the logic of sarcasm:} \\
\texttt{    *   Sarcasm often relies on **subverting expectations** or **using a known phrase in a way that highlights absurdity**.} \\
\texttt{    *   If the prompt implies that "feeding trolls" is a common, annoying task, saying you are "glad" to do it is already sarcastic.} \\
\texttt{    *   However, the specific twist in these types of logic puzzles often lies in the **absurd substitution**.} \\
\texttt{    *   Let's look closer at the "Trial \#1" reflection provided in the prompt: *"I misidentified the sarcastic target; in option (A), the phrase 'feed the elephants' is a nonsensical or absurd substitution... making (A) the intended answer."*} \\
\texttt{    *   The prompt explicitly states in the "Trial \#1" section that the correct answer was **(A)** and explains *why*: the absurdity of substituting "elephants" for "trolls" creates the humor/sarcasm by treating a nonsensical action as a standard social interaction, or by highlighting the absurdity of the situation through a deliberate misuse of the idiom. The "endangered" part applies to elephants (which are actually endangered), making the statement a mix of literal truth (elephants are endangered) and absurd context (feeding them as if they were internet trolls), which creates a specific type of satirical humor.} \\
\texttt{    *   Conversely, Option (} \\
            }
        }
    }
    & & \\ \\

    \theutterance \stepcounter{utterance}  
    & & & \multicolumn{2}{p{0.3\linewidth}}{
        \cellcolor[rgb]{0.9,0.9,0.9}{
            \makecell[{{p{\linewidth}}}]{
                \texttt{\tiny{[GM$|$GM]}}
                \texttt{Target: (B)} \\
            }
        }
    }
    & & \\ \\

    \theutterance \stepcounter{utterance}  
    & & & \multicolumn{2}{p{0.3\linewidth}}{
        \cellcolor[rgb]{0.9,0.9,0.9}{
            \makecell[{{p{\linewidth}}}]{
                \texttt{\tiny{[GM$|$GM]}}
                \texttt{game\_result = LOSE} \\
            }
        }
    }
    & & \\ \\

    \theutterance \stepcounter{utterance}  
    & & \multicolumn{4}{p{0.6\linewidth}}{
        \cellcolor[rgb]{0.95,0.95,0.95}{
            \makecell[{{p{\linewidth}}}]{
                \texttt{\tiny{[P1$\leftrightarrow$P]}}
                \texttt{Post{-}episode Reflexion (not sent to GM):} \\
\texttt{New plan: My previous failure stemmed from over{-}analyzing the "endangered" descriptor in option (B) and missing the primary sarcastic mechanism: the absurd substitution of the idiom itself. In the provided examples, sarcasm is consistently generated by stating a falsehood that mocks a situation (e.g., claiming ad hominem attacks help, or comparing league punishments to politics). In the final question, option (A) replaces the standard idiom "feed the trolls" with "feed the elephants." This creates a nonsensical, absurd image that mocks the act of engaging with online provocateurs by treating it as a literal, impossible, or ridiculous task (feeding elephants), whereas option (B) uses the correct idiom with a secondary ironic adjective. Next time, I will prioritize identifying the option that fundamentally breaks the expected idiom or logic to create absurdity, rather than looking for subtle ironic adjectives within a correct idiom. I will specifically check if one option swaps a key noun in a common phrase to create a humorous, illogical statement, as this is the stronger indicator of the intended sarcasm in this specific dataset.} \\
            }
        }
    }
    & \\ \\

\end{supertabular}
}

\end{document}
